\DeclareEditionOverlay{CM-C03-U001}{%
\begin{EditionUnitCard}{De secante a tangente}
Sigue cómo cambia la pendiente cuando el segundo punto se acerca al primero.
\RegisteredBookFigure{FIG-C03-001}{fig:c03:secant-tangent}{De la pendiente secante a la pendiente tangente.}{Dos puntos de una curva determinan una secante; cuando el incremento h tiende a cero, la secante se aproxima a la tangente en el punto base.}{\SecantTangentVisual}
\begin{EditionWarningBox}{Atención}
h se aproxima a cero; el cociente incremental no se evalúa en h igual a cero.
\end{EditionWarningBox}
\end{EditionUnitCard}}

\DeclareEditionOverlay{CM-C03-U002}{%
\begin{EditionUnitCard}{Una lupa lineal}
Cerca del punto, la recta tangente organiza el cambio principal.
\begin{EditionCheckpointBox}{Comprueba}
Identifica punto base, incremento y término lineal antes de aproximar.
\end{EditionCheckpointBox}
\end{EditionUnitCard}}

\DeclareEditionOverlay{CM-C03-U003}{%
\begin{EditionUnitCard}{Una implicación, no una equivalencia}
La derivabilidad garantiza continuidad, pero una función continua puede tener esquina.
\begin{EditionWarningBox}{Atención}
No inviertas la implicación sin una hipótesis adicional.
\end{EditionWarningBox}
\end{EditionUnitCard}}

\DeclareEditionOverlay{CM-C03-U004}{%
\begin{EditionUnitCard}{Patrones con condiciones}
Relaciona cada fórmula con una derivación modelo y con su dominio.
\begin{EditionCheckpointBox}{Comprueba}
Antes de usar una regla elemental, verifica dónde está definida la función.
\end{EditionCheckpointBox}
\end{EditionUnitCard}}

\DeclareEditionOverlay{CM-C03-U005}{%
\begin{EditionUnitCard}{Reglas con hipótesis}
Las fórmulas son herramientas condicionadas, no identidades sin contexto.
\begin{EditionCheckpointBox}{Comprueba}
Nombra qué funciones deben ser derivables y qué denominador debe ser no nulo.
\end{EditionCheckpointBox}
\end{EditionUnitCard}}

\DeclareEditionOverlay{CM-C03-U006}{%
\begin{EditionUnitCard}{Derivar otra vez también exige existencia}
Cada orden añade una condición de regularidad.

\end{EditionUnitCard}}

\DeclareEditionOverlay{CM-C03-U007}{%
\begin{EditionUnitCard}{Desmonta la composición}
Nombra las capas antes de derivar de fuera hacia dentro.
\begin{EditionCheckpointBox}{Comprueba}
Marca la función interior y verifica que aparece su derivada.
\end{EditionCheckpointBox}
\end{EditionUnitCard}}

\DeclareEditionOverlay{CM-C03-U008}{%
\begin{EditionUnitCard}{Inversa no significa recíproco}
Sigue qué punto corresponde en la función original y revisa las hipótesis locales.
\begin{EditionWarningBox}{Atención}
Distingue $f^{-1}$ de $1/f$ y de la preimagen de un conjunto.
\end{EditionWarningBox}
\end{EditionUnitCard}}

\DeclareEditionOverlay{CM-C03-U009}{%
\begin{EditionUnitCard}{Una relación puede definir localmente}
Deriva la identidad y conserva la condición que permite aislar la derivada.
\begin{EditionCheckpointBox}{Comprueba}
Identifica el factor por el que se divide y comprueba que no sea cero.
\end{EditionCheckpointBox}
\end{EditionUnitCard}}

\DeclareEditionOverlay{CM-C03-U010}{%
\begin{EditionUnitCard}{Transformar con dominio}
La derivación logarítmica simplifica la estructura, pero no elimina las restricciones.
\begin{EditionWarningBox}{Atención}
Escribe primero dónde el logaritmo natural tiene sentido.
\end{EditionWarningBox}
\end{EditionUnitCard}}

\DeclareEditionOverlay{CM-C03-U011}{%
\begin{EditionUnitCard}{Un cero produce un candidato}
La condición de Fermat ayuda a localizar; no clasifica por sí sola.
\begin{EditionWarningBox}{Atención}
f'(a)=0 no basta para afirmar máximo o mínimo.
\end{EditionWarningBox}
\end{EditionUnitCard}}

\DeclareEditionOverlay{CM-C03-U012}{%
\begin{EditionUnitCard}{Tres hipótesis, una tangente horizontal}
Comprueba el intervalo y los extremos antes de buscar el punto interior.
\begin{EditionCheckpointBox}{Comprueba}
Explica qué puede fallar si eliminas cada hipótesis.
\end{EditionCheckpointBox}
\end{EditionUnitCard}}

\DeclareEditionOverlay{CM-C03-U013}{%
\begin{EditionUnitCard}{Cambio medio y cambio instantáneo}
El teorema garantiza algún punto interior con la pendiente adecuada.
\begin{EditionCheckpointBox}{Comprueba}
Antes de aplicarlo, verifica continuidad y derivabilidad en sus intervalos correctos.
\end{EditionCheckpointBox}
\end{EditionUnitCard}}

\DeclareEditionOverlay{CM-C03-U014}{%
\begin{EditionUnitCard}{l'Hôpital es un teorema controlado}
Reconoce la forma, comprueba hipótesis y vuelve a evaluar después de derivar.
\begin{EditionWarningBox}{Atención}
No se aplica porque haya un cociente, sino porque se cumplen sus condiciones.
\end{EditionWarningBox}
\end{EditionUnitCard}}

\DeclareEditionOverlay{CM-C03-U015}{%
\begin{EditionUnitCard}{Revisa después de cada paso}
Una aplicación cambia el problema; vuelve a diagnosticarlo.
\begin{EditionCheckpointBox}{Comprueba}
Indica la forma obtenida antes de decidir si repites el método.
\end{EditionCheckpointBox}
\end{EditionUnitCard}}

\DeclareEditionOverlay{CM-C03-U016}{%
\begin{EditionUnitCard}{Primero transforma, después decide}
Las formas producto y diferencia no entran directamente en la regla.
\begin{EditionWarningBox}{Atención}
Justifica la transformación y revisa la nueva forma.
\end{EditionWarningBox}
\end{EditionUnitCard}}

\DeclareEditionOverlay{CM-C03-U017}{%
\begin{EditionUnitCard}{Exponenciales indeterminadas}
Trabaja primero con el logaritmo del objeto cuando el dominio lo permite.

\end{EditionUnitCard}}

